% ======================================================================
% R01 -- Git
% Molde: por que -> modelo -> nucleo -> exemplo -> armadilhas
%        -> exercicios -> checklist
% ======================================================================
\modulo[R01 Git]{r01}{R01 -- Git}

% ----------------------------------------------------------------------
\subsection[Por que importa]{Por que isso importa aqui}
% ----------------------------------------------------------------------
\begin{frame}{Por que Git nesta disciplina}
  \begin{itemize}
    \item Tudo o que produzimos é \alert{texto}: VHDL, Tcl, MATLAB, Python, \LaTeX.
    \item Precisamos \alert{voltar} ao estado em que o testbench passava,
      \alert{comparar} versões e \alert{experimentar} sem estragar o que funciona.
    \item O repositório é o \alert{caderno de laboratório}: cada commit registra
      o quê, quando e \emph{por quê}.
    \item A CI verifica cada mudança; as releases entregam o PDF destas notas.
  \end{itemize}
  \vfill
  \begin{block}{Analogia}
    O Git é um álbum de fotos de uma obra: você arruma o enquadramento
    (\emph{preparação}), tira a foto (\emph{commit}) e ela nunca mais muda.
  \end{block}
\end{frame}

% ----------------------------------------------------------------------
\subsection[Modelo mental]{Modelo mental: snapshots e três áreas}
% ----------------------------------------------------------------------
\begin{frame}{Snapshots, não diferenças}
  \centering
  \begin{tikzpicture}[
      c/.style={circle, draw=labsemMarinho, fill=labsemGelo,
        minimum size=9mm, font=\scriptsize},
      >=Stealth]
    \node[c] (a) {a1f};
    \node[c, right=12mm of a] (b) {7c2};
    \node[c, right=12mm of b] (d) {e90};
    \draw[<-, thick] (a) -- (b);
    \draw[<-, thick] (b) -- (d);
    \node[draw, rounded corners, fill=labsemCobre!20, above=5mm of d,
      font=\scriptsize] (m) {main};
    \node[draw, rounded corners, fill=labsemVerde!20, right=4mm of m,
      font=\scriptsize] (h) {HEAD};
    \draw[->] (m) -- (d);
    \draw[->] (h) -- (m);
  \end{tikzpicture}
  \vspace{4mm}
  \begin{itemize}
    \item Cada \alert{commit} é uma fotografia do projeto inteiro.
    \item Cada commit aponta para o \alert{pai}: a história é um grafo.
    \item Identificador = \alert{hash SHA} do conteúdo. Mudou um byte, muda o hash.
    \item Branch = \alert{etiqueta móvel} que aponta para um commit;
      \cmd{HEAD} = onde você está.
  \end{itemize}
\end{frame}

\begin{frame}{As três áreas}
  \centering
  \begin{tikzpicture}[
      area/.style={draw=labsemMarinho, fill=labsemGelo, rounded corners,
        minimum width=30mm, minimum height=13mm, align=center, font=\small},
      >=Stealth]
    \node[area] (wd) {Diretório de\\trabalho};
    \node[area, right=15mm of wd] (ix) {Preparação\\(\emph{index})};
    \node[area, right=15mm of ix] (rp) {Repositório\\(\cmd{.git})};
    \draw[->, thick] (wd) -- node[above, font=\scriptsize] {\cmd{git add}} (ix);
    \draw[->, thick] (ix) -- node[above, font=\scriptsize] {\cmd{git commit}} (rp);
    \draw[->, thick, labsemCobre] (rp.south) to[bend left=25]
      node[below, font=\scriptsize] {\cmd{git restore} / \cmd{git switch}} (wd.south);
  \end{tikzpicture}
  \vspace{3mm}
  \begin{itemize}
    \item \textbf{Trabalho}: os arquivos no disco, com suas edições.
    \item \textbf{Preparação}: a \alert{próxima foto} sendo montada;
      você escolhe o que entra.
    \item \textbf{Repositório}: fotos já tiradas, imutáveis.
  \end{itemize}
  \begin{alertblock}{Consequência}
    \cmd{git commit} fotografa a \emph{preparação}, não o disco.
  \end{alertblock}
\end{frame}

% ----------------------------------------------------------------------
\subsection[Ciclo e commits]{Ciclo básico e commits}
% ----------------------------------------------------------------------
\begin{frame}[fragile]{O ciclo que se repete sempre}
  \begin{columns}[T]
    \begin{column}{0.56\textwidth}
      \begin{lstlisting}[style=bash]
        git status            # onde está cada arquivo
        git diff              # trabalho -> preparação
        git add -p arq        # prepara trecho a trecho
        git diff --staged     # preparação -> commit
        git commit -m "..."   # tira a foto
        git lg                # log --oneline --graph
      \end{lstlisting}
    \end{column}
    \begin{column}{0.40\textwidth}
      \begin{block}{Commit atômico}
        Um commit = um assunto. Corrigiu o testbench e ajustou um
        comentário no RTL? \alert{Dois commits}.
      \end{block}
      \small
      Por quê: quando algo quebrar, acha-se o commit culpado e ele é
      desfeito sem levar junto o que estava certo.
    \end{column}
  \end{columns}
\end{frame}

\begin{frame}[fragile]{Mensagens: Conventional Commits}
  \begin{lstlisting}[style=bash]
    tipo(escopo): descrição no imperativo   (<= 72 caracteres)

    corpo opcional: explica o PORQUÊ
  \end{lstlisting}
  \small
  \begin{tabularx}{\textwidth}{@{}lX@{}}
    \toprule
    \cmd{feat} / \cmd{fix}               & funcionalidade / correção \\
    \cmd{docs} / \cmd{test}              & documentação e slides / testbenches \\
    \cmd{refactor} / \cmd{style}         & reorganização / formatação \\
    \cmd{build} / \cmd{ci} / \cmd{chore} & toolchain / pipelines / manutenção \\
    \bottomrule
  \end{tabularx}
  \vspace{2mm}

  \textcolor{labsemVerde}{\ding{51}}~\cmd{fix(tcl): compila corpo do pacote após a declaração}\\
  \textcolor{labsemVermelho}{\ding{55}}~\cmd{ajustes}
  \vfill
  \footnotesize
  Escopos: \cmd{vhdl tcl matlab python slides docs ci repo r01 m2 proj1}.
  Quebra de compatibilidade: \cmd{feat(tcl)!:}.
\end{frame}

% ----------------------------------------------------------------------
\subsection[Higiene]{Higiene: ignorar e normalizar}
% ----------------------------------------------------------------------
\begin{frame}{\cmd{.gitignore}, \cmd{.gitattributes}, \cmd{.editorconfig}}
  \begin{tabularx}{\textwidth}{@{}lX@{}}
    \toprule
    \cmd{.gitignore}     & o que as ferramentas \alert{geram} fica fora:
                           \cmd{db/}, \cmd{output\_files/}, \cmd{work/}, \cmd{*.wlf},
                           \cmd{transcript}, \cmd{*.asv}, \cmd{\_\_pycache\_\_/},
                           \cmd{*.aux}, \cmd{out/}\ldots \\
    \cmd{.gitattributes} & finais de linha: \alert{LF} para tudo, CRLF só em
                           \cmd{.bat/.cmd}; binários marcados;
                           \cmd{.do/.sdc/.qsf} contados como Tcl \\
    \cmd{.editorconfig}  & o editor segue as mesmas regras (UTF-8, LF, recuo
                           por linguagem) \\
    \bottomrule
  \end{tabularx}
  \vspace{3mm}
  \begin{block}{Regra}
    O que é \alert{gerado} não se versiona: reconstrói-se com um comando.
    Git LFS só se um dia for preciso versionar dado bruto grande (ADR).
  \end{block}
\end{frame}

% ----------------------------------------------------------------------
\subsection[Branches e merge]{Branches e merge}
% ----------------------------------------------------------------------
\begin{frame}[fragile]{Branches: trabalhar em paralelo}
  \begin{columns}[T]
    \begin{column}{0.5\textwidth}
      \centering
      \begin{tikzpicture}[
          c/.style={circle, draw=labsemMarinho, fill=labsemGelo, minimum size=6mm},
          f/.style={circle, draw=labsemVerde, fill=labsemVerde!15, minimum size=6mm},
          >=Stealth]
        \node[c] (a) {};
        \node[c, right=6mm of a] (b) {};
        \node[f, above right=4mm and 6mm of b] (f1) {};
        \node[f, right=6mm of f1] (f2) {};
        \node[c, below right=4mm and 22mm of b] (m) {};
        \draw[<-] (a) -- (b);
        \draw[<-] (b) -- (f1);
        \draw[<-] (f1) -- (f2);
        \draw[<-] (b) -- (m);
        \draw[<-] (f2) -- (m);
        \node[font=\scriptsize, right=1mm of m] {merge};
        \node[font=\scriptsize, above=0mm of f2, text=labsemVerde] {feat/x};
      \end{tikzpicture}
    \end{column}
    \begin{column}{0.48\textwidth}
      \begin{lstlisting}[style=bash]
        git switch -c feat/x   # cria e entra
        # ... commits ...
        git switch main
        git merge feat/x
        git branch -d feat/x
      \end{lstlisting}
    \end{column}
  \end{columns}
  \begin{itemize}
    \small
    \item \alert{Conflito}: os dois lados mudaram as mesmas linhas.
      Edite o arquivo, apague \cmd{<<<<<<< ======= >>>>>>>},
      \cmd{git add}, \cmd{git commit}.
    \item No GitHub usamos \alert{squash merge}: o branch vira um único
      commit no \cmd{main}.
  \end{itemize}
\end{frame}

% ----------------------------------------------------------------------
\subsection[Remoto e fluxo]{Remoto e fluxo profissional}
% ----------------------------------------------------------------------
\begin{frame}{Remoto e fluxo de trabalho}
  \begin{columns}[T]
    \begin{column}{0.45\textwidth}
      \small
      \begin{tabularx}{\linewidth}{@{}lX@{}}
        \toprule
        \cmd{fetch} & baixa, não mexe nos arquivos \\
        \cmd{pull}  & fetch + avança (\cmd{pull.ff only}) \\
        \cmd{push}  & envia commits \\
        \bottomrule
      \end{tabularx}
      \vspace{2mm}

      Autenticação por \alert{SSH} (\cmd{gh auth login}).
    \end{column}
    \begin{column}{0.52\textwidth}
      \centering
      \begin{tikzpicture}[
          s/.style={draw=labsemMarinho, fill=labsemGelo, rounded corners,
            font=\scriptsize, minimum width=18mm, minimum height=6mm},
          >=Stealth, node distance=3mm]
        \node[s] (i) {issue};
        \node[s, below=of i] (b) {branch};
        \node[s, below=of b] (p) {PR};
        \node[s, right=6mm of p] (c) {CI verde};
        \node[s, above=of c] (sq) {squash merge};
        \node[s, above=of sq] (t) {tag \(\rightarrow\) release};
        \draw[->] (i) -- (b);
        \draw[->] (b) -- (p);
        \draw[->] (p) -- (c);
        \draw[->] (c) -- (sq);
        \draw[->] (sq) -- (t);
      \end{tikzpicture}
    \end{column}
  \end{columns}
  \vspace{2mm}
  \small
  \cmd{main} sempre verde, protegido por \emph{ruleset}. Versões \alert{SemVer}:
  \cmd{0.MINOR.PATCH} durante o curso; \cmd{1.0.0} = projeto final.
\end{frame}

% ----------------------------------------------------------------------
\subsection[Hooks e CI/CD]{Automação: hooks e CI/CD}
% ----------------------------------------------------------------------
\begin{frame}{Verificação automática em três níveis}
  \begin{tabularx}{\textwidth}{@{}lXl@{}}
    \toprule
    Nível                  & O que verifica                                   & Quando \\
    \midrule
    \cmd{pre-commit}       & espaços, fim de arquivo, LF, YAML, arquivos
                             grandes, chaves privadas, mensagem convencional  & a cada commit \\
    CI (\cmd{ci.yml})      & os mesmos hooks + (cresce) testes Python,
                             regressão VHDL, compilação do deck               & a cada PR/push \\
    CD (\cmd{release.yml}) & release com notas automáticas + PDF das notas    & a cada tag \cmd{v*} \\
    \bottomrule
  \end{tabularx}
  \vspace{3mm}
  \begin{block}{Por que a CI cresce aos poucos}
    Cada job entra junto com o conteúdo que ele verifica: você entende cada
    linha do YAML e o \cmd{main} nunca fica vermelho por padrão.
  \end{block}
\end{frame}

% ----------------------------------------------------------------------
\subsection[Desfazer]{Desfazer com segurança}
% ----------------------------------------------------------------------
\begin{frame}{Desfazer com segurança}
  \small
  \begin{tabularx}{\textwidth}{@{}Xll@{}}
    \toprule
    Situação                                   & Comando                                     & Destrói? \\
    \midrule
    Descartar edição não preparada             & \cmd{git restore arq}                       & a edição \\
    Tirar da preparação                        & \cmd{git restore --staged arq}              & não \\
    Corrigir o último commit (não enviado)     & \cmd{git commit --amend}                    & reescreve \\
    Desfazer N commits locais, mantendo mudanças & \cmd{git reset --soft HEAD\textasciitilde N} & não \\
    Desfazer commit já enviado                 & \cmd{git revert <hash>}                     & não \\
    Guardar trabalho pela metade               & \cmd{git stash} / \cmd{pop}                 & não \\
    ``Perdi um commit''                        & \cmd{git reflog}                            & não \\
    \bottomrule
  \end{tabularx}
  \vspace{2mm}
  \begin{alertblock}{Regra}
    Nunca reescrever (\cmd{amend}, \cmd{reset}, \emph{force push}) o que já
    está no \cmd{main} remoto.
  \end{alertblock}
\end{frame}

% ----------------------------------------------------------------------
\subsection[Armadilhas]{Armadilhas}
% ----------------------------------------------------------------------
\begin{frame}{Armadilhas}
  \begin{armadilha}
    \begin{itemize}
      \small
      \item \cmd{git commit -a} pula a preparação: entra tudo o que é
        rastreado e mudou.
      \item Caminhos com espaço/acentos quebram Quartus e ModelSim no Windows.
      \item \cmd{core.autocrlf} ligado + \cmd{.gitattributes}: finais de linha
        trocados a cada commit. Deixe o \cmd{.gitattributes} mandar.
      \item Versionar \cmd{work/}, \cmd{db/} ou PDFs: repositório inchado e
        conflitos inúteis.
      \item Mensagem padrão do \cmd{git revert} (\cmd{Revert "..."}) é recusada
        pelo hook: reescreva como \cmd{revert: ...}.
      \item Ruleset em repositório privado exige GitHub Pro (Student Pack).
    \end{itemize}
  \end{armadilha}
\end{frame}

% ----------------------------------------------------------------------
\subsection[Exercícios]{Exercícios e checklist}
% ----------------------------------------------------------------------
\begin{frame}{Exercícios}
  \begin{enumerate}
    \small
    \item Você editou um arquivo depois do último commit e não rodou
      \cmd{git add}. O que \cmd{git commit} registra?
      \irpara{r01-resp1}{Resposta}
    \item Após \cmd{git add}, você editou o mesmo arquivo de novo. O que aparece
      em \cmd{git diff}, em \cmd{git diff --staged} e o que entra no commit?
      \irpara{r01-resp2}{Resposta}
    \item Provoque um conflito entre dois branches num repositório de rascunho
      e resolva-o. Mostre o \cmd{git lg}.
    \item Um commit ruim já está no \cmd{main} remoto. Qual comando usar e por quê?
      \irpara{r01-resp4}{Resposta}
  \end{enumerate}
\end{frame}

\begin{frame}{Checklist}
  \begin{seisei}
    \begin{itemize}
      \small
      \item explico as três áreas e o que cada comando move entre elas;
      \item faço commits atômicos com mensagens convencionais;
      \item crio branch, abro PR, espero a CI e faço squash merge pelo \cmd{gh};
      \item sei o que o \cmd{.gitignore} e o \cmd{.gitattributes} do repositório fazem;
      \item resolvo um conflito sem pânico;
      \item escolho entre \cmd{restore}, \cmd{reset}, \cmd{revert} e \cmd{stash};
      \item publico uma release por tag.
    \end{itemize}
  \end{seisei}
\end{frame}
